\scanpage{35}
\uline{Lemme 1}
\[
\left\{\begin{aligned}
&G\ \text{\deleted{\unsure{séparé}, \illegible{2} connexe}\inserted{\uline{irred} et \uline{propre}} sur $k$,}\\
&e\ \text{pt \unsure{ratln|ratl} de $G$}\\
&G\times G\longrightarrow G,\quad\text{loi de multiplication \deleted{ayant} $e$ comme unité}\\
&G\xrightarrow{\vee}G^{\deleted{\unsure{x}}}\quad\text{application \deleted{\unsure{inverse bilatère} à droite}}\\
&\hspace{12em}\inserted{\text{telle que }\ xy^{\vee}=e\iff x=y}
\end{aligned}\right.
\]
\uline{Alors} \deleted{\unsure{$G$|c'est}} la loi de composition est associative et commutative [donc $G$ est une variété abélienne]

\uline{Démonstration} Considérons
\[
\deleted{f(a,b,c)=[(ab)c][a(bc)]^{\vee}}
\]
\[
f(a,b)=(ab)(ba)^{\vee}
\]
on a
\[
f(e,\unsure{x})=f(x,e)=e
\]
donc
\[
\left\{\begin{aligned}
&f(a,x)=f(e,x)=e\quad\text{pour }a\in U\ \text{voisinage de }e\\
&[f(x,a)=f(x,e)=e\qquad\text{---}\quad a\in U'\quad\text{---------}\,]
\end{aligned}\right.
\]
\margin{inutile}
\deleted{donc $f(x,x)=f(e,x)=f(x,e)=f(e,e)=$}

donc par \uline{\unsure{prolongmt}} [fonction rationnelle] on aura
\[
f(x,y)=e\qquad\text{pour tt }x,y\quad\text{d'où}
\]
donc $(ab)(ba)^{\vee}=e$ i.e. $ba=ab$\hfill commutativité

Pour l'associativité, on fait
\[
f(a,b,c)=[(ab)c][a(bc)]^{\vee}
\]
on aura
\[
f(a,b,c)=e\qquad\text{si }\unsure{a,[b\ \text{ou}\ c]|a,[\text{ou } b\ \text{ou }c]}=e
\]
donc
\[
f(\unsure{x|a},\unsure{y|b},a)=e\qquad
\text{\deleted{pour fonction \unsure{rat.}}\inserted{pour $a\in U$ vois. de $e$}}
\]
donc
\[
f(abc)=
\]
% archive sheet number: 34
